\documentclass[11pt]{article}

% ---------- Packages ----------
\usepackage[margin=1in]{geometry}
\usepackage{amsmath, amssymb, amsthm}
\usepackage{mathtools}
\usepackage{enumitem}
\usepackage{graphicx}
\usepackage{booktabs}
\usepackage{array}
\usepackage{hyperref}
\usepackage{xcolor}
\usepackage{fancyhdr}
\usepackage{lastpage}
\usepackage{listings}
\usepackage{courier}

% ---------- Hyperref ----------
\hypersetup{
    colorlinks=true,
    linkcolor=blue,
    urlcolor=blue
}

% ---------- Header/Footer ----------
\pagestyle{fancy}
\fancyhf{}

\lhead{MIT 6.1200J}
\chead{Problem Set 9}
\rhead{Nicolas Iskandar}

\cfoot{\thepage\ /\ \pageref{LastPage}}

% ---------- Custom Commands ----------
\newcommand{\R}{\mathbb{R}}
\newcommand{\N}{\mathbb{N}}
\newcommand{\Z}{\mathbb{Z}}

\newcommand{\answer}{
    \vspace{0.5em}
    \noindent\textbf{Answer.}
}

% ---------- Code Style ----------
\lstset{
    basicstyle=\ttfamily\small,
    breaklines=true,
    frame=single,
    columns=fullflexible,
    keepspaces=true
}

% ---------- Title ----------
\title{
    \vspace{-2em}
    MIT 6.1200J \\
    Problem Set 9
}
\author{Nicolas Iskandar}
\date{}

% ==================================================
\begin{document}

\maketitle

\section*{Problem 1: Double Counting}

We must prove that for all natural numbers $n \le a \le b \le c$:
$$ \sum_{i,j,k} \binom{a}{i} \binom{b}{j} \binom{c}{k} = \binom{a+b+c}{n} $$
where the sum ranges over all $i, j, k \in \mathbb{N}$ satisfying $i+j+k=n$.

\textbf{Combinatorial Proof:}
Suppose we have a set of $a+b+c$ distinct items. Let this set be partitioned into three disjoint subsets: Group A of size $a$, Group B of size $b$, and Group C of size $c$. 

We want to count the total number of ways to choose exactly $n$ items from this entire set of $a+b+c$ items.

\textbf{Method 1 (Right-Hand Side):} 
By definition, the number of ways to choose $n$ items from a set of $a+b+c$ items is simply $\binom{a+b+c}{n}$. 

\textbf{Method 2 (Left-Hand Side):}
Alternatively, we can categorize our choices based on how many items we take from each of the three groups. Let $i$ be the number of items chosen from Group A, $j$ be the number of items chosen from Group B, and $k$ be the number of items chosen from Group C. To choose exactly $n$ items in total, it must be true that $i + j + k = n$.

For any fixed valid triplet $(i, j, k)$:
\begin{itemize}
    \item We can choose $i$ items from Group A in $\binom{a}{i}$ ways.
    \item We can choose $j$ items from Group B in $\binom{b}{j}$ ways.
    \item We can choose $k$ items from Group C in $\binom{c}{k}$ ways.
\end{itemize}
By the Rule of Product, there are exactly $\binom{a}{i} \binom{b}{j} \binom{c}{k}$ ways to select the $n$ items such that the group breakdown is exactly $i, j,$ and $k$. 

To find the total number of ways to choose the $n$ items across all possible valid distributions, we sum over all combinations of $i, j,$ and $k$ that sum to $n$. This gives the expression:
$$ \sum_{i,j,k} \binom{a}{i} \binom{b}{j} \binom{c}{k} $$

Since both methods count the exact same combinatorial set, the two expressions must be equal.

\vspace{0.5cm}

\section*{Problem 2: Pigeonhole Principle}

We are given a subset $S \subseteq \{1, 2, 3, \dots, 2n\}$ where $|S| = n + 1$. We want to show that there exist distinct $a, b \in S$ such that $a/b$ is a power of 2.

Following the hint to factor each number into an odd number and a power of 2, any integer $x \in S$ can be uniquely expressed in the form $x = q \cdot 2^k$, where $q$ is an odd integer and $k \ge 0$ is a non-negative integer.

We structure our Pigeonhole Principle argument as follows:
\begin{itemize}
    \item \textbf{The Pigeons:} The $n+1$ distinct elements of the set $S$.
    \item \textbf{The Holes:} The possible odd factors for numbers in the range $1$ to $2n$. These are the odd integers $1, 3, 5, \dots, 2n-1$. Notice that there are exactly $n$ odd numbers in this range.
    \item \textbf{The Rule:} Assign each pigeon $x \in S$ to the hole corresponding to its odd factor $q$ (i.e., if $x = q \cdot 2^k$, it goes into hole $q$).
\end{itemize}

Since we are placing $n+1$ pigeons into $n$ holes, by the Pigeonhole Principle, at least one hole must contain at least two pigeons. 

This means there are two distinct elements $a, b \in S$ that share the exact same odd factor $q$. We can write them as:
$$ a = q \cdot 2^i $$
$$ b = q \cdot 2^j $$
Since $a$ and $b$ are distinct elements of $S$, $i \neq j$. Without loss of generality, assume $i > j$. 
Then, dividing $a$ by $b$ gives:
$$ \frac{a}{b} = \frac{q \cdot 2^i}{q \cdot 2^j} = 2^{i-j} $$
Since $i > j$, the difference $i-j$ is a positive integer, meaning $a/b$ is a power of 2. This completes the proof.

\vspace{0.5cm}

\section*{Problem 3: Dice Dice Dice Dice Dice}

\begin{enumerate}[label=\textbf{(\alph*)}]
    \item \textbf{Probability Space:}
    \begin{itemize}
        \item \textbf{Outcomes:} The set of outcomes consists of all ordered 5-tuples $(d_1, d_2, d_3, d_4, d_5)$ representing the results of the 5 dice rolls, where each $d_i \in \{1, 2, 3, 4, 5, 6\}$.
        \item \textbf{Number of outcomes:} Because there are 5 independent dice and each has 6 possible faces, there are $6^5 = 7776$ possible outcomes.
        \item \textbf{Probability:} Assuming the standard six-sided dice are fair, each outcome is equally likely. Thus, the probability of any single specific outcome is uniform: $\frac{1}{6^5} = \frac{1}{7776}$.
    \end{itemize}

    \item $A_{29}$ is the event that the sum of the dice is at least 29. The maximum possible sum is $5 \times 6 = 30$. Thus, $A_{29}$ includes all outcomes where the sum is exactly 29 or 30.
    \begin{itemize}
        \item Sum = 30 requires all sixes: $(6, 6, 6, 6, 6)$.
        \item Sum = 29 requires four sixes and one five. There are 5 such permutations based on which die is the 5.
    \end{itemize}
    Explicitly, $A_{29} = \{(6,6,6,6,6), (5,6,6,6,6), (6,5,6,6,6), (6,6,5,6,6), (6,6,6,5,6), (6,6,6,6,5)\}$.
    Since there are exactly 6 outcomes in $A_{29}$, we have:
    $$ \Pr[A_{29}] = \frac{6}{6^5} = \frac{1}{6^4} = \frac{1}{1296} $$

    \item $A_{18}$ is the event that the sum of the dice is at least 18. 
    Let $A_{<18}$ be its complement (the event that the sum is exactly 17 or less).
    We define a function $f: A_{18} \to A_{<18}$ mapping an outcome $(d_1, d_2, d_3, d_4, d_5)$ to $(7-d_1, 7-d_2, 7-d_3, 7-d_4, 7-d_5)$. 
    
    \textbf{Proof it is a bijection:} 
    First, verify it maps into the target set. If the sum of the original dice is $S \ge 18$, the sum of the new mapped dice is $\sum_{i=1}^5 (7 - d_i) = 35 - S$. Since $S \ge 18$, the new sum $35 - S \le 35 - 18 = 17$. Thus, the output is strictly in $A_{<18}$. 
    Second, verify it is invertible. Applying the function twice yields $7 - (7 - d_i) = d_i$. Because it is its own inverse, it maps 1-to-1 and onto, making it a bijection.

    Because a bijection exists between $A_{18}$ and its complement, both sets must have the exact same number of elements. Thus, $|A_{18}| = \frac{1}{2} |S_{total}|$, which means:
    $$ \Pr[A_{18}] = \frac{1}{2} $$

    \item Let $B_k$ be the event of exactly $k$ sixes, and $C_k$ be the event of at least $k$ sixes. 
    To find $B_k$, we choose which $k$ of the 5 dice show a six ($\binom{5}{k}$ ways). The chosen $k$ dice have a $1/6$ chance of being six, and the remaining $5-k$ dice have a $5/6$ chance of not being six. Thus:
    $$ \Pr[B_k] = \binom{5}{k} \left(\frac{1}{6}\right)^k \left(\frac{5}{6}\right)^{5-k} = \frac{\binom{5}{k} 5^{5-k}}{6^5} $$
    Since $C_k$ is the cumulative probability of having $k$ or more sixes:
    $$ \Pr[C_k] = \sum_{i=k}^5 \Pr[B_i] = \sum_{i=k}^5 \frac{\binom{5}{i} 5^{5-i}}{6^5} $$

    \item We want $\Pr[C_3 \mid C_2]$. Using the definition of conditional probability and noting that $C_3 \subset C_2$ (if you have at least 3 sixes, you definitely have at least 2), $C_3 \cap C_2 = C_3$.
    $$ \Pr[C_3 \mid C_2] = \frac{\Pr[C_3 \cap C_2]}{\Pr[C_2]} = \frac{\Pr[C_3]}{\Pr[C_2]} $$
    Using the verified values given in the problem:
    $$ \Pr[C_3 \mid C_2] = \frac{276/6^5}{1526/6^5} = \frac{276}{1526} = \frac{138}{763} \approx 0.1809 $$

    \item Let $D_k$ be the event that Ben rolls at least $k$ dice showing the same value. We want $\Pr[D_3 \mid D_2]$. 
    First, find the total size of $D_3$. There are 6 values a die can show. For each value, the number of ways to get at least 3 of them is equal to $|C_3| = 276$. Since it is impossible to have at least 3 of one value AND at least 3 of a different value at the same time (which would require 6 dice), these 6 events are disjoint. 
    Thus, $|D_3| = 6 \times 276 = 1656$.
    
    Next, find $|D_2|$. It is easier to find its complement, $D_{<2}$, which means no two dice show the same value (all 5 dice are distinct). 
    $|D_{<2}| = 6 \times 5 \times 4 \times 3 \times 2 = 720$. 
    Thus, $|D_2| = 6^5 - 720 = 7776 - 720 = 7056$.

    Since $D_3 \subset D_2$:
    $$ \Pr[D_3 \mid D_2] = \frac{|D_3 \cap D_2|}{|D_2|} = \frac{|D_3|}{|D_2|} = \frac{1656}{7056} = \frac{23}{98} \approx 0.2347 $$

    \item Ben is down to 3 dice. The new sample space size is $6^3 = 216$. The pairs of dice faces that sum exactly to 6 are $\{1,5\}$, $\{2,4\}$, and $\{3,3\}$. 
    Let $E_{15}$ be the event that the 3 rolls include a 1 and a 5.
    Let $E_{24}$ be the event that the 3 rolls include a 2 and a 4.
    Let $E_{33}$ be the event that the 3 rolls include at least two 3s.
    We seek the size of the union $|E_{15} \cup E_{24} \cup E_{33}|$. 
    
    Using Inclusion-Exclusion:
    \begin{itemize}
        \item $|E_{15}|$: We can have (1,5,X). There are 3! = 6 permutations if the 3 dice are distinct, and 4 choices for X (not 1 or 5), giving $6 \times 4 = 24$. If X is 1 (e.g., 1,1,5), there are 3 permutations. If X is 5, there are 3 permutations. Total $|E_{15}| = 24 + 3 + 3 = 30$.
        \item $|E_{24}|$: By identical symmetry to $\{1,5\}$, $|E_{24}| = 30$.
        \item $|E_{33}|$: We can have (3,3,X) where X is not 3. There are $\binom{3}{2} = 3$ ways to place the 3s, and 5 choices for X. $3 \times 5 = 15$. Plus the single outcome of (3,3,3) makes 16. Total $|E_{33}| = 16$.
    \end{itemize}
    Next, look at intersections. Is it possible to satisfy two of these events simultaneously with only 3 dice? 
    $E_{15} \cap E_{24}$ requires a 1, 5, 2, and 4 (4 dice). Thus, it's empty (size 0).
    $E_{15} \cap E_{33}$ requires a 1, 5, 3, 3. Empty (size 0).
    $E_{24} \cap E_{33}$ requires 2, 4, 3, 3. Empty (size 0).
    
    The total number of valid outcomes is simply $30 + 30 + 16 = 76$. 
    $$ \Pr[\text{pair sums to 6}] = \frac{76}{216} = \frac{19}{54} $$
\end{enumerate}

\vspace{0.5cm}

\section*{Problem 4: Unreliable Helpers}

\begin{enumerate}[label=\textbf{(\alph*)}]
    \item We translate the problem statement into formal probability equations. 
    Given $\Pr[E] = 0.10$ and $\Pr[\bar{E}] = 0.90$. 
    
    The TeaBot malfunctions 20\% of the time. This means it gives the opposite answer of reality:
    $$ \Pr[T \mid \bar{E}] = 0.20 \quad \text{(Says error when there isn't one)} $$
    $$ \Pr[\bar{T} \mid E] = 0.20 \implies \Pr[T \mid E] = 1 - 0.20 = 0.80 \quad \text{(Says error when there is one)} $$
    
    The LecBot malfunctions 25\% of the time:
    $$ \Pr[L \mid \bar{E}] = 0.25 $$
    $$ \Pr[\bar{L} \mid E] = 0.25 \implies \Pr[L \mid E] = 1 - 0.25 = 0.75 $$

    \item We want to find $\Pr[E \mid \bar{T} \cap L]$. 
    Since the bots malfunction independently given the state of the problem, we can use conditional independence: 
    $\Pr[\bar{T} \cap L \mid E] = \Pr[\bar{T} \mid E] \Pr[L \mid E] = (0.20)(0.75) = 0.15$
    $\Pr[\bar{T} \cap L \mid \bar{E}] = \Pr[\bar{T} \mid \bar{E}] \Pr[L \mid \bar{E}] = (0.80)(0.25) = 0.20$

    By Bayes' Theorem:
    $$ \Pr[E \mid \bar{T} \cap L] = \frac{\Pr[\bar{T} \cap L \mid E] \Pr[E]}{\Pr[\bar{T} \cap L \mid E] \Pr[E] + \Pr[\bar{T} \cap L \mid \bar{E}] \Pr[\bar{E}]} $$
    $$ \Pr[E \mid \bar{T} \cap L] = \frac{(0.15)(0.10)}{(0.15)(0.10) + (0.20)(0.90)} $$
    $$ \Pr[E \mid \bar{T} \cap L] = \frac{0.015}{0.015 + 0.180} = \frac{0.015}{0.195} = \frac{15}{195} = \frac{1}{13} $$

    \item Is $\Pr[T \mid L]$ equal to $\Pr[T]$?
    
    \textbf{Intuition:} No, they are not equal. While the robots malfunction independently, they are both evaluating the exact same underlying problem. If the LecBot says there is an error ($L$), this is evidence that the problem likely does indeed have an error ($E$). Because the problem is now more likely to have an error, the TeaBot (who is mostly accurate) is correspondingly more likely to also report an error. Thus, observing $L$ provides information about $E$, which in turn provides information about $T$, breaking unconditional independence. We expect $\Pr[T \mid L] > \Pr[T]$.
    
    \textbf{Numerical Verification:}
    First, find the unconditional probability $\Pr[T]$ using the Law of Total Probability:
    $$ \Pr[T] = \Pr[T \mid E]\Pr[E] + \Pr[T \mid \bar{E}]\Pr[\bar{E}] $$
    $$ \Pr[T] = (0.80)(0.10) + (0.20)(0.90) = 0.08 + 0.18 = 0.26 $$
    
    Next, we find $\Pr[L]$:
    $$ \Pr[L] = \Pr[L \mid E]\Pr[E] + \Pr[L \mid \bar{E}]\Pr[\bar{E}] $$
    $$ \Pr[L] = (0.75)(0.10) + (0.25)(0.90) = 0.075 + 0.225 = 0.30 $$
    
    Next, we find the intersection $\Pr[T \cap L]$ using conditional independence:
    $$ \Pr[T \cap L] = \Pr[T \cap L \mid E]\Pr[E] + \Pr[T \cap L \mid \bar{E}]\Pr[\bar{E}] $$
    $$ \Pr[T \cap L] = (\Pr[T \mid E]\Pr[L \mid E])\Pr[E] + (\Pr[T \mid \bar{E}]\Pr[L \mid \bar{E}])\Pr[\bar{E}] $$
    $$ \Pr[T \cap L] = (0.80 \times 0.75)(0.10) + (0.20 \times 0.25)(0.90) $$
    $$ \Pr[T \cap L] = (0.60)(0.10) + (0.05)(0.90) = 0.06 + 0.045 = 0.105 $$
    
    Now we calculate $\Pr[T \mid L]$:
    $$ \Pr[T \mid L] = \frac{\Pr[T \cap L]}{\Pr[L]} = \frac{0.105}{0.30} = 0.35 $$
    
    Since $0.35 \neq 0.26$, we mathematically verify that $\Pr[T \mid L] \neq \Pr[T]$.
\end{enumerate}

\end{document}